%!TEX TS-program = xelatex
\documentclass[]{friggeri-cv}
\usepackage{enumitem}
\usepackage{hyperref}
\usepackage{color}
\usepackage{xcolor}
\hypersetup{
	pdftitle={},
	pdfauthor={},
	pdfsubject={},
	pdfkeywords={},
	colorlinks=false,       % no lik border color
	allbordercolors=white    % white border color for all
}
\addbibresource{bibliography.bib}
\RequirePackage{xcolor}
\definecolor{pblue}{HTML}{0395DE}

\begin{document}
\header{Chao}{Li}
{Software Engineer}

% Fake text to add separator      
\fcolorbox{white}{gray}{\parbox{\dimexpr\textwidth-2\fboxsep-2\fboxrule}{%
	.....
}}

% In the aside, each new line forces a line break
\begin{aside}
	\section{Education}
	\textbf{Master of Science in 
	Computer Science}, 
	2013-2015
	Worcester Polytechnic 
	Institute, US
	~
	\textbf{Bachelor in 
	Computer Science}, 
	2008-2012
	Tianjin University of 
	Technology, China
	\section{Skill set}
	C++\\ JAVA\\ Python\\ MapReduce\\ CUDA\\ SQL\\ Sockets programming\\ Relational databases\\ etc.
	~
	\section{OS Preference}
	\textbf{Windows}\includegraphics[scale=0.40]{img/5stars.png}
	\textbf{Unix/Linux}\includegraphics[scale=0.40]{img/4stars.png}
	\textbf{MacOS}\includegraphics[scale=0.40]{img/4stars.png}
	~
	\section{Address}
	1302 N Capitol AVE,
	Unit 1,
	San Jose, 95132
	\section{Tel}
	+1(774)329-5650
	\section{Mail}
	\href{mailto:lichaorodxx@gmail.com}{\textbf{lichaorodxx@}\\gmail.com}
	~
	\section{Web}
	\href{http://users.wpi.edu/~cli5}{users/wpi.edu/\textasciitilde cli5}
	\href{https://www.linkedin.com/in/lichao90}{linkedin.com/in/lichao90}
	\includegraphics[scale=0.05]{img/Chao_Li.png}
\end{aside}

\section{Projects}
\begin{entrylist}
	\entry
	{02/14 - 05/14}
	{Network protocol simulation and application}
	{Computer networks course project}
	{
		\vspace{-3mm}
		\begin{itemize}
			\item Implemented a server and a client of a chat app, based on Linux
				sockets;
			\item Simulated the \emph{selective repeat} and \emph{go back n}
				algorithms of TCP protocal, running on a emulated linked layer.
			\item This program is multithread, which can handle multiple clients at
				the same time.
		\end{itemize}
	}
	\entry
	{09/10 - 10/11}
	{Study on the Object Tracking in Video Stream}
	{College student innovation project}
	{
		\vspace{-3mm}
		\begin{itemize}
			\item Studied Meanshift algorithm and Particle Filter algorithm;
			\item Implemented both of the above two algorithms in C++ and OpenCV;
			\item Blended the two algorithms.
		\end{itemize}
	}
\end{entrylist}
\section{Research experience}
\begin{entrylist}
	\entry
	{09/14 - 05/15}
	{Semidefinite Programming, Binary Codes and a Graph Coloring Problem}
	{Master's thesis project Advised by Prof. Martin, William J}
	{
		\vspace{-3mm}
		\begin{itemize}
			\item Studied {\em Semidefinite Programming} (SDP), and CVX used to model
				and solve SDPs;
			\item Studied the binary code size bound formulation based on SDP;
			\item Wrote a C++ program used to generate the dual-SDP problems
				automatically.\\
		\end{itemize}
	}
	\entry
	{06/14 - 02/15}
	{Recovering All Statistically Significant OPSM Patterns}
	{Advised by Prof. Trapp, Andrew}
	{
		\vspace{-3mm}
		\begin{itemize}
			\item Studied the Order-Preserved Submatrix (OPSM) problem;
			\item Used C++ and CPLEX Concert technology to formulate its complementary problem;
			\item Added constraints to recover all statistically significant OPSMs.\\
		\end{itemize}
	}
	\entry
	{06/14 - 12/14}
	{Study on the Combinatorics of Binary Serieal Parallel Graph}
	{Advised by Prof. Hofri, Micha}
	{
		\vspace{-3mm}
		\begin{itemize}
			\item Used the \emph{Symbolic method} in {\em Analytic Combinatorics} to
				model random structures;
			\item Used the \emph{Transfer Theorem} in {\em Analytic Combinatorics} to
				obtain the analytical generating functions for the structure;
			\item Studied on the generating function for its properties, e.g.,size and
				order.\\
		\end{itemize}
	}
\end{entrylist}
\begin{entrylist}
	\entry
	{04/13 - 05/13}
	{Internship}
	{Tellyes Company, Research Engineer, Tianjin, China}
	{
		\vspace{-3mm}
		\begin{itemize}
			\item Learned the convolutional based ultrasound image simulation
				algorithm;
			\item Rewrote the algorithm from C++ into CUDA C giving a 10 times faster
				improvement.
		\end{itemize}
	}
\end{entrylist}
\section{Publications}
Micha Hofri, Chao Li, Hosam Mahmoud\\
\textbf{On the combinatorics of binary serires-parallel graphs}\\
\emph{Submitted to Probability in the Engineering and Informational Science.}


\end{document}
